\section{The Tangible Immediacy of the Present}

\begin{quotex}
This is final section of an article by \textbf{Guido de Giorgio}, titled The Instant and Eternity, first published in Diorama Filosofico in 1939.
\end{quotex}


How many of Dante's admirers are there who are not content in glorifying his verses or expression---something absolutely exterior and superficial---but who in applying the doctrine, the knowledge on all the planes of being to which they relate and in the totality of the Heavenly Voyage.

The past is nothing if it is not integrated, lived, validated by personal experience, by life, if it is not totalized and exalted in the great shaking of the eternal now. The moderns, when they are not fornicating in the past like the thieves in a necropolis, turn their backs on it, contemplating then the hypothetical ``sun of the future'' which will never shine, because the future exists only as the last evanescent border of an arduous vision, a mirage and nothing more, a fallacious projection stained by the spasms of their own insufficiency. The ``non-achievement'' in the face of the Truth, the incurable sentiment of one who does not know nor wants to know, does not know how to nor wants to carry with himself all the weight of the world, to assume it in the divine instant, created the myth of the future. Turning obstinately their backs to that which is, they wait with curiosity for what is not, for what will be, and they long for the confirmation of a dream by an illusory reflection of the dream itself, in a nocturnal march of phantoms that alone engender the present, by the spontaneity of its flux and its mirage. Strange speculation on the future, which makes them forget the treasures of the past and the tangible immediacy of the present. For they are really only in the present, with all worlds, in the essential unity of the point, jewel of all jewels, the eternal eye of God!

We would like to still say some other things, but we prefer to conclude with the words of \textbf{Zarathustra}:

\begin{quotex}
To these men of today will I not be light, nor be called light. Them, will I blind with the lightning of my wisdom! Put out their eyes!
\flright{\textsc{Friedrich Nietzsche}, \emph{Zarathustra}}
\end{quotex}



\flrightit{Posted on 2012-06-12 by Cologero}

\begin{center}* * *\end{center}

\begin{footnotesize}
\begin{sffamily}
\texttt{Rajan Meru on 2026-02-11 at 20:25 said:}

Evocative and prescient. A really excellent essay.

Hope to see more of De Giorgio's works translated into English.

\end{sffamily}
\end{footnotesize}

